% ============================================
% Johns Hopkins Letter Template
% Assistant Professor Cover Letter – DePauw University
% ============================================

\documentclass[11pt,letterpaper]{letter}

\usepackage{graphicx}
\usepackage{eso-pic}
\usepackage{geometry}
\usepackage{microtype}
\usepackage[hidelinks]{hyperref}

% ----- Page Setup -----
\geometry{left=25mm,right=25mm,top=25mm,bottom=25mm}

% ----- Logo Placement (first page only) -----
\newcommand{\PutJHULogoFG}[1]{%
  \AddToShipoutPictureFG*{%
    \AtPageUpperLeft{%
      \hspace{10mm}%
      \raisebox{-38mm}{%
        \includegraphics[width=6cm]{#1}%
      }%
    }%
  }%
}

% ----- Sender Information -----
\signature{Decory Edwards}
\address{Department of Economics \\ Johns Hopkins University \\ Baltimore, MD 21218 \\ \texttt{dedwar65@jh.edu}}

\begin{document}
\PutJHULogoFG{Images/jhulogo.png}

\begin{letter}{Search Committee \\
Department of Economics and Management \\
DePauw University \\
Greencastle, IN 46135 \\ USA}

\opening{Dear Members of the Search Committee,}

I am writing to express my interest in the tenure-track Assistant Professor position in the Department of Economics and Management at DePauw University, beginning August 2026. I am a Ph.D. candidate in Economics at Johns Hopkins University, advised by Professors Christopher D.~Carroll, Nicholas W.~Papageorge, and Stelios Fourakis, and I expect to receive my degree in May 2026. My research fields are macroeconomics, wealth and income dynamics, and computational economics.

My research agenda focuses on understanding the determinants of wealth inequality and the mechanisms through which heterogeneity in returns to wealth shapes aggregate outcomes. In my job-market paper, I estimate the degree of heterogeneity in individual portfolio returns using micro-data from the Survey of Consumer Finances and simulate a structural model in which households face idiosyncratic return risk. I show that allowing for persistent heterogeneity in returns substantially improves the model’s ability to match the empirical distribution of wealth, offering a tractable way to connect micro-level return dispersion to macroeconomic inequality.

In a second project, I use the Health and Retirement Study to examine the relationship between interpersonal trust and portfolio performance. Building on the literature that finds a hump-shaped relationship between trust and income, I show that a similar relationship holds for individual returns to wealth measured in the HRS data, highlighting a behavioral channel through which social attitudes shape saving and investment outcomes. This work also provides new estimates of the persistent component of returns to net worth, which motivate the calibration of heterogeneity in my structural model.

I am particularly enthusiastic about the opportunity to join DePauw University’s School of Business and Leadership, where rigorous undergraduate instruction in finance is complemented by a vibrant liberal-arts environment and a commitment to experiential and interdisciplinary learning. I welcome the opportunity to contribute to courses such as financial economics, financial crises, financial econometrics, behavioral finance, corporate finance, portfolio management, and fintech, as well as core offerings including \textit{Investment Analysis and Portfolio Management}, \textit{Introduction to Economics}, \textit{Introduction to Econometrics}, \textit{Intermediate Microeconomic Theory}, \textit{Intermediate Macroeconomic Theory}, and the \textit{Senior Seminar}. I am also eager to support undergraduate advising, develop electives that integrate data-driven financial applications, and contribute to DePauw’s mission of preparing students for meaningful leadership in business and society.

\textbf{Teaching.} My teaching experience centers on undergraduate macroeconomics and an upper-level macro policy seminar, where I have led weekly discussion sections and held regular office hours for classes ranging from 16 to 40 students. My approach is student-centered: I aim to help students “convince themselves” that they have moved from confusion to understanding by holding structured office-hour coaching that builds trust and confidence. At DePauw, I am prepared to teach \textit{Financial Economics}, \textit{Financial Econometrics}, \textit{Behavioral Finance}, \textit{Corporate Finance}, \textit{Portfolio Management}, and related courses in addition to the economics and finance core. I welcome opportunities to supervise undergraduate research, integrate real-world data into coursework, and contribute to the department’s applied learning initiatives.

I believe I would be a strong fit because my computational and empirical experience allows me to (i) teach finance and economics using real-world datasets and reproducible quantitative methods, (ii) design active learning environments that combine theory with applied decision-making, and (iii) contribute meaningfully to a liberal-arts community that values collaboration, diverse perspectives, and experiential learning. I am especially drawn to DePauw’s commitment to diversity, inclusive pedagogy, and the development of ethical and capable leaders through the liberal arts.

I would be honored to join the Department of Economics and Management at DePauw University. Thank you for your consideration; I welcome the opportunity to discuss my fit with your department at your convenience.

\closing{Sincerely,}

\end{letter}
\end{document}
